\chapter{Spesifikasi Kebutuhan Sistem}
\label{app:spesifikasi}

Lampiran ini berisi spesifikasi kebutuhan fungsional dan nonfungsional dari sistem \emph{chatbot} hukum yang dikembangkan. Spesifikasi disusun berdasarkan modul utama sistem yang saling terpisah namun terintegrasi dalam satu \emph{workflow} \emph{end-to-end}.

\section{Spesifikasi Kebutuhan Fungsional Sistem}

Berikut adalah spesifikasi kebutuhan fungsional sistem berdasarkan modul utama yang telah didefinisikan.

{\setlength{\tabcolsep}{4pt}
\begin{longtable}{|p{0.09\textwidth}|p{0.34\textwidth}|p{0.49\textwidth}|}
\caption{Spesifikasi Kebutuhan Fungsional Sistem \emph{Chatbot} Hukum}\label{tab:functional}\endfirsthead
\hline \textbf{ID} & \textbf{Modul} & \textbf{Kebutuhan} \\\hline
\endhead
\hline
F01 & Modul data \emph{ingestion} (\emph{scraping} dan \emph{parsing}) serta modul \emph{admin} dan \emph{landing page} & Modul dapat mengumpulkan dokumen peraturan perundang-undangan dari sumber resmi. \\\hline
F02 & Modul data \emph{ingestion} (\emph{scraping} dan \emph{parsing}) serta modul \emph{admin} dan \emph{landing page} & Modul dapat melakukan \emph{cleaning} dan normalisasi dokumen peraturan perundang-undangan. \\\hline
F03 & Modul data \emph{ingestion} (\emph{scraping} dan \emph{parsing}) serta modul \emph{admin} dan \emph{landing page} & Modul dapat melakukan \emph{parsing} struktur hukum secara hierarkis menjadi undang-undang, bab, bagian, pasal, dan ayat. \\\hline
F04 & Modul \emph{Document Processing Orchestrator} (DPO) dan \emph{query} pre-planning & Modul dapat melakukan segmentasi dokumen hukum menjadi unit makna yang lebih kecil. \\\hline
F05 & Modul \emph{Document Processing Orchestrator} (DPO) dan \emph{query} pre-planning & Modul dapat melakukan \emph{embedding} terhadap unit dokumen hukum beserta \emph{metadata} strukturnya. \\\hline
F06 & Modul \emph{Document Processing Orchestrator} (DPO) dan \emph{query} pre-planning & Modul dapat melakukan \emph{guardrails} dan sanitasi PII. \\\hline
F07 & Modul \emph{Document Processing Orchestrator} (DPO) dan \emph{query} pre-planning & Modul dapat melakukan deteksi \emph{intent} dan transformasi \emph{query} agar selaras dengan struktur hukum formal. \\\hline
F08 & Modul \emph{hybrid retrieval} dan \emph{Knowledge Graph} serta modul autentikasi & Modul dapat melakukan \emph{hybrid retrieval} untuk menemukan \emph{chunk} pasal yang relevan dengan \emph{query}. \\\hline
F09 & Modul \emph{hybrid retrieval} dan \emph{Knowledge Graph} serta modul autentikasi & Modul dapat melakukan penelusuran relasi hukum melalui \emph{Knowledge Graph}. \\\hline
F10 & Modul \emph{hybrid retrieval} dan \emph{Knowledge Graph} serta modul autentikasi & Modul dapat memperluas konteks hukum berdasarkan relasi yang telah didapatkan pada \emph{Knowledge Graph}. \\\hline
F11 & Lapisan \emph{Explainable AI} (XAI) serta modul \emph{chat} & Modul dapat menghasilkan jawaban hukum berdasarkan konteks hasil \emph{retrieval} dan \emph{reasoning}. \\\hline
F12 & Lapisan \emph{Explainable AI} (XAI) serta modul \emph{chat} & Modul dapat menampilkan dasar hukum berupa undang-undang, pasal, dan ayat yang dirujuk. \\\hline
F13 & Lapisan \emph{Explainable AI} (XAI) serta modul \emph{chat} & Modul dapat menyediakan mekanisme \emph{explainability} berupa \emph{source attribution} dan jalur \emph{reasoning}. \\\hline
F14 & Lapisan Explainable AI (XAI) serta modul \emph{chat} & Modul dapat menyajikan penjelasan tambahan berbasis \emph{input saliency} dan \emph{counterfactual}. \\\hline
\end{longtable}}

\section{Spesifikasi Kebutuhan Nonfungsional Sistem}

Berikut adalah spesifikasi kebutuhan nonfungsional sistem \emph{chatbot} hukum.

\begin{table}[H]
\centering
\caption{Spesifikasi Kebutuhan Nonfungsional Sistem}
\label{tab:nonfunctional}
\begin{tabularx}{\textwidth}{|p{0.14\textwidth}|p{0.25\textwidth}|X|}
\hline \textbf{ID} & \textbf{Parameter} & \textbf{Kebutuhan} \\\hline
NF01 & \emph{Usability} & Sistem menyediakan antarmuka yang dapat digunakan oleh pengguna umum. \\\hline
NF02 & \emph{Usability} & Sistem menampilkan pesan kesalahan dan informasi sistem dalam Bahasa Indonesia yang jelas. \\\hline
NF03 & \emph{Reliability} & Sistem tetap dapat memberikan umpan balik kepada pengguna ketika terjadi kegagalan sebagian modul. \\\hline
NF04 & \emph{Reliability} & \emph{Auto-save draft query} setiap 5 detik. \\\hline
\end{tabularx}
\end{table}
